% 2010 年湖南卷理科第 13 题：体积 20 cm^3 的几何体的三视图。
% 来源 PDF：普通高考/2010/2010湖南理.pdf，SHA-256 db6274a80ebb0022ff1b91b2cd3714fe8cca335e48be69805dcd87cdefd8c8d1。
% 原页：第 1 页（A3 三栏）第 2 栏下部，书页码 601。
% 重绘依据：三个直角三角形——正视图直角边 5 与 h（直角在左下），
%   侧视图直角边 6 与 h（直角在左下），俯视图直角边 5、6（直角在左上）；
%   俯视图在原页未标数字，此处亦不标。图中 h 按原页比例虚设，不代表答案。
\begin{tikzpicture}[scale=0.62, semithick]
  % 正视图：宽 5，高 h
  \draw (0,0) -- (5,0) -- (0,3.2) -- cycle;
  \draw[<->] (-0.6,0) -- (-0.6,3.2);
  \draw (-0.75,0) -- (-0.35,0);
  \draw (-0.75,3.2) -- (-0.35,3.2);
  \node[left] at (-0.85,1.6) {$h$};
  \draw[<->] (0,-0.6) -- (5,-0.6);
  \draw (0,-0.75) -- (0,-0.35);
  \draw (5,-0.75) -- (5,-0.35);
  \node[below] at (2.5,-0.55) {$5$};
  \node[below] at (2.5,-1.45) {正视图};

  % 侧视图：宽 6，高 h
  \draw (7.6,0) -- (13.6,0) -- (7.6,3.2) -- cycle;
  \draw[<->] (7.6,-0.6) -- (13.6,-0.6);
  \draw (7.6,-0.75) -- (7.6,-0.35);
  \draw (13.6,-0.75) -- (13.6,-0.35);
  \node[below] at (10.6,-0.55) {$6$};
  \node[below] at (10.6,-1.45) {侧视图};

  % 俯视图：宽 5，高 6，直角在左上
  \draw (0,-4.6) -- (5,-4.6) -- (0,-10.6) -- cycle;
  \node[below] at (2.5,-10.9) {俯视图};
\end{tikzpicture}
